\subsection{Wave-specific concentration of supercritical human transmission}
\label{sec:supp_wave_amplification}

To examine whether periods of inferred supercritical human transmission were
preferentially associated with epidemic expansion, we divided the mortality
series into three major epidemic waves and classified each wave into an
ascending and a declining phase. Wave boundaries were defined by the
inter-wave mortality troughs, and the observed mortality maximum within each
wave was used to separate ascent from decline. We then summarized the
posterior median human reproduction number, $R_{0,H}$, together with the
inferred human and zoonotic forces of infection within each phase
(Table~\ref{tab:supp_wave_amplification}).

The clearest phase-specific separation was observed in the second and third
epidemic waves. During wave~2, the posterior median of $R_{0,H}$ exceeded
one in 51.4\% of ascending-phase weeks, compared with
7.1\% of declining-phase weeks. Using the more
stringent criterion that the lower 2.5\% posterior quantile also exceeded one,
the corresponding proportions were 34.3\% and
0.0\%, respectively. A similar pattern was
observed in wave~3, for which 45.0\% of ascending
weeks had a posterior median $R_{0,H}>1$, whereas no declining-phase week
did so. Under the stricter posterior criterion, the corresponding proportions
were 30.0\% and
0.0\%.

Wave~1 showed a weaker separation between phases. The posterior median
$R_{0,H}$ exceeded one in 25.0\% of ascending
weeks and 17.4\% of declining weeks, while the
stricter criterion was met in 15.0\% and
13.0\%, respectively. Thus, the concentration
of supercritical human transmission during epidemic ascent was most evident
for the second and third waves rather than being equally pronounced across
the full study period.

The timing of inferred human transmission was also consistent with a role for
human amplification in the later epidemic waves. Human FOI reached its
within-wave maximum 0 weeks
at the wave~2 mortality peak
and 4 weeks
before the wave~3 mortality peak.
For wave~1, the human FOI maximum occurred
4 weeks
after the mortality peak.
Importantly, inferred human FOI did not collapse immediately after the
mortality peak. This is not inconsistent with $R_{0,H}<1$: the force of
infection depends on the number of infectious humans present, so substantial
human transmission can persist while existing transmission chains are
contracting.

The inferred spillover component showed a different temporal pattern.
Within waves~2 and~3, the maximum spillover FOI occurred
3 and
2 weeks after the mortality
peak, respectively, whereas in wave~1 it occurred
7 weeks before the mortality
peak. The later waves therefore do not support a simple interpretation in
which a single large spillover pulse immediately precedes each mortality
wave. Instead, conditional on the fitted model, they are consistent with
continued zoonotic introductions occurring alongside a temporally varying
capacity for human-to-human amplification.

Taken together, these wave-specific summaries suggest that the second and
third mortality waves were characterized by a marked concentration of
supercritical human transmission during their ascending phases, followed by
predominantly subcritical human transmission after the mortality peak despite
continued inferred spillover. This pattern is consistent with recurrent
zoonotic seeding combined with episodic human amplification. Because all
transmission quantities in this analysis are posterior estimates from the
same fitted mechanistic model, however, the wave analysis should be
interpreted as a model-conditional description of the inferred mechanism
rather than as an independent causal validation.
